%! Zeros \& Graphs of Polynomials — Unit 3, Lesson 3.6 — Extra Credit (Answer Key)
% MIRROR of ../extra_credit/main.tex: identical except -boxes -> -key and the page-header title.
% Every \gradedblank line, work block and box size is copied VERBATIM; key mode prints the answers.
% UNITS 3+ PACKET-FIRST shape (LESSON_SHAPE.md §0a, 2026-11-14). GENERATED, GRADED, OPTIONAL.
% Last in the packet; the build starts it on a fresh sheet so it tears off. TWO PAGES, exactly,
% blank and key. Five stretch items, harder than every ★★★:
%   q1     COMBINE  Rounds 1+2: a double real zero and an imaginary zero → degree-4 standard form
%   q2a/b  REVERSE  Rounds 3+1: equation of a pre-drawn quartic from its graph; its imaginary count
%   q3     SOL-style MULTIPLE CHOICE on the Fundamental Theorem and conjugate pairs
%   q4a-c  MODEL    a cross-country course profile; interpret h(4) and the below-start stretch
%   q5     FIND THE MISTAKE  a touch treated as multiplicity 1 (rubric)
% Every final answer is a \gradedblank, the line copied VERBATIM into the key (§9). Work blocks
% sit OUTSIDE the boxes. NO NAME LINE. Contexts differ from the notes' (teach, then transfer).
\documentclass[12pt]{article}
\usepackage{algebra2-article}
\usepackage{algebra2-key}
\usepackage{grading}          % math-shared: key mode prints each answer
\gradingassignment{alg2-u03-l06-ec}
\usepackage{algebra2-graded}
\gradingcourse{algebra_2}
% Pre-drawn polynomial graph (identical in the key) --- the notes' macros.
\newcommand{\pgrid}[5]{%
  \draw[linegray!45, very thin, xstep=1, ystep=#5] (#1,#3) grid (#2,#4);
  \draw[->, semithick, charcoal] (#1,0)--(#2,0) node[right, font=\tiny, inner sep=1pt]{$x$};
  \draw[->, semithick, charcoal] (0,#3)--(0,#4) node[above, font=\tiny, inner sep=1pt]{$y$};}
\newcommand{\pcurve}[5]{%
  \begin{scope}\clip (#1,#3) rectangle (#2,#4);
    \draw[very thick, forest, domain=#1:#2, samples=160, smooth] plot (\x,{#5});
  \end{scope}}
\newcommand{\xtick}[1]{\foreach \t in {#1}{%
  \draw[charcoal] (\t,0) +(0,-2pt) -- +(0,2pt);
  \node[below, font=\tiny, inner sep=1.5pt, fill=white, fill opacity=0.8, text opacity=1]
    at (\t,0) {$\t$};}}
\newcommand{\ytick}[1]{\foreach \t in {#1}{%
  \draw[charcoal] (0,\t) +(-2pt,0) -- +(2pt,0);
  \node[left, font=\tiny, inner sep=1.5pt, fill=white, fill opacity=0.8, text opacity=1]
    at (0,\t) {$\t$};}}

\begin{document}

\pageheader{Unit 3, Lesson 3.6}{Extra Credit --- Answer Key\strut}

\vspace{0.10in}

\boxguard[8]
\begin{remindbox}
\small
\textbf{Extra credit --- optional.} Tear off these pages and turn them in with your homework,
the first class after two study halls. I will announce the due date in class and on TurtleNet.
{\small\itshape Write each final answer in its box. Only what is inside a box is graded.}\par
\end{remindbox}

\vspace{0.06in}

\begin{notesbox}{Stretch}
\small
Every polynomial on these pages has real coefficients. Show your work outside the boxes.
\begin{enumerate}[leftmargin=1.9em, itemsep=8pt]
  \item \textbf{Count, then build.} A polynomial $p$ has leading coefficient $1$ and degree $4$.
    Two of its zeros are $3$, a \emph{double} zero, and $1-2i$. Write $p(x)$ in standard form.
    \begin{work}
      p(x) &= (x-3)^2(x-1+2i)(x-1-2i) \\
           &= (x^2-6x+9)(x^2-2x+5) \\
           &= x^4-8x^3+26x^2-48x+45
    \end{work}
    $p(x)=$ \inblank{\blank{7.0cm}}{\gradedblank[id=q1, type=short, points=2, answer={x^4 - 8x^3 + 26x^2 - 48x + 45}, check=expression]{7.0cm-2pt}{12.5pt}}

  \item \textbf{Read it backward.} The graph shows a polynomial of degree $4$. Its
    $y$-intercept is $-4$.\par
    \noindent
    \begin{minipage}[t]{0.60\linewidth}
    \begin{work}
      f(x) &= a(x+1)^2(x-2)(x-4) \\
      -4 &= a(1)^2(-2)(-4) \\
      -4 &= 8a \\
      a &= -\tfrac12
    \end{work}
    \end{minipage}\hfill
    \begin{minipage}[t]{0.36\linewidth}
    \vspace{0pt}\centering
    \begin{tikzpicture}[x=0.36cm, y=0.085cm]
      \pgrid{-2}{5}{-14}{10}{2}
      \pcurve{-2}{5}{-14}{10}{-0.5*((\x)+1)*((\x)+1)*((\x)-2)*((\x)-4)}
      \xtick{-1,2,4} \ytick{-4}
    \end{tikzpicture}
    \end{minipage}\par
    \vspace{4pt}
    (a) Write its equation in factored form. \quad $f(x)=$ \inblank{\blank{5.6cm}}{\gradedblank[id=q2a, type=short, points=2, answer={-1/2 (x + 1)^2 (x - 2)(x - 4)}, check=expression]{5.6cm-2pt}{12.5pt}}

    \vspace{6pt}
    (b) How many imaginary zeros does $f$ have? \inblank{\blank{1.6cm}}{\gradedblank[id=q2b, type=short, points=1, answer={0}, check=numeric]{1.6cm-2pt}{12.5pt}}

  \item \textbf{Multiple choice (SOL-style).} A polynomial $q$ has real coefficients and
    degree $4$. Which statement \emph{must} be true?
    Write the letter of the single best answer in its Letter box.
    \begin{enumerate}[label=(\Alph*), leftmargin=2.0em, itemsep=1pt]
      \item $q$ has exactly four real zeros.
      \item If $2+3i$ is a zero of $q$, then $2-3i$ is also a zero of $q$.
      \item The graph of $q$ crosses the $x$-axis at least once.
      \item $q$ can have exactly three imaginary zeros.
    \end{enumerate}
    \letterbox{\gradedblank[id=q3, type=choice, points=1, answer={B}]{0.4in}{0.4in}}
\end{enumerate}
\end{notesbox}

\newpage
\begin{notesbox}{Stretch, continued}
\small
\begin{enumerate}[leftmargin=1.9em, itemsep=8pt, start=4]
  \item \textbf{Model it.} The cross-country team's new course climbs and dips. Its height
    above the starting line is $h(x)$ meters, $x$ hundred meters into the course
    ($0\le x\le6$). The course is at starting-line height only at $x=0$, $2$, and $6$, and
    $h$ is a cubic. At $x=1$ the course is $5$ meters above the start.
    \begin{work}
      h(x) &= a\,x(x-2)(x-6) \\
      5 &= a(1)(-1)(-5) \\
      a &= 1 \\
      h(4) &= 4(2)(-2)=-16
    \end{work}
    (a) $h(x)=$ \inblank{\blank{4.4cm}}{\gradedblank[id=q4a, type=short, points=2, answer={x(x - 2)(x - 6)}, check=expression]{4.4cm-2pt}{12.5pt}}
    \qquad (b) $h(4)=$ \inblank{\blank{2.0cm}}{\gradedblank[id=q4b, type=short, points=1, answer={-16}, check=numeric]{2.0cm-2pt}{12.5pt}}

    \vspace{6pt}
    (c) Explain what $h(4)$ means on the course, and say for which stretch of the course the
    runners are below the starting line.
    \rubricbox{\gradedblank[id=q4c, type=rubric, points=2, display={$h(4)=-16$ means that $400$ meters into the course the runners are $16$ meters \emph{below} the starting line. The cubic is negative between its zeros $2$ and $6$ (it crosses at both), so the course is below the start from $200$ m to $600$ m.}]{\linewidth}{2in}}

  \item \textbf{Find the mistake.} A graph of a cubic \emph{touches} the $x$-axis at $x=-2$,
    \emph{crosses} it at $x=3$, and passes through $(0,-12)$. Student A wrote:
    \begin{center}
    \emph{``The zeros are $-2$ and $3$, so $f(x)=a(x+2)(x-3)$. Then $-12=a(2)(-3)=-6a$, so
    $a=2$ and $f(x)=2(x+2)(x-3)$.''}
    \end{center}
    Which step is wrong? Explain why, and write the correct equation.
    \rubricbox{\gradedblank[id=q5, type=rubric, points=3, display={The first step is wrong: Student A treated the touch at $x=-2$ as a single zero. A touch means an even multiplicity, so the factor is $(x+2)^2$; a cubic also needs three zeros counted with multiplicity, and $-2,-2,3$ gives three. Then $-12=a(2)^2(-3)=-12a$, so $a=1$ and $f(x)=(x+2)^2(x-3)$.}]{\linewidth}{2in}}
\end{enumerate}
\end{notesbox}

\end{document}
